\labsection{7}{Link aggregation: several cables, one link}{sec:lab07-lacp}

\begin{labproject}{Bundle physical links, then make the bundle negotiate}
\textbf{Coverage.} Chapter~\ref{ch:lacp}.\\
\textbf{Materials.} \texttt{Lab07/manual\_mode.txt}, \texttt{Lab07/lacp\_mode.txt}.\\
\textbf{Goal.} Build an Eth-Trunk in manual mode, rebuild it in LACP mode, and understand why the second is safer.

\medskip\noindent\textbf{Part A --- Two problems, one solution.}\par
Aggregation solves bandwidth and redundancy together. Four gigabit links behave as one four-gigabit link, and losing one degrades capacity instead of connectivity.

\begin{keyidea}
Aggregation also solves a problem you would otherwise create. Connect two switches with four parallel cables and no aggregation, and you have built four loops --- spanning tree will block three of them, leaving you exactly where you started with three cables wasted.

An Eth-Trunk presents the bundle to spanning tree as a single logical interface, so all four links forward.
\end{keyidea}

\medskip\noindent\textbf{Part B --- Traffic does not split evenly.}\par

\begin{pitfall}
A four-link trunk does not give one flow four times the bandwidth.

The switch hashes header fields --- source and destination MAC, IP, or port depending on configuration --- and every frame in a given conversation lands on the same physical link. This is deliberate: splitting a flow across links would reorder packets and hurt TCP badly.

So a single large file transfer uses one link. Aggregation increases \emph{aggregate} capacity across many flows, not the speed of any one of them. Expect uneven utilisation, and choose a hash that matches your traffic --- MAC-based hashing is useless behind a router, where every frame carries the same router MAC.
\end{pitfall}

\medskip\noindent\textbf{Part C --- Manual mode.}\par

\begin{lstlisting}[style=dcnterminal]
system-view
interface Eth-Trunk 1
 mode manual load-balance
quit

interface GigabitEthernet 0/0/1
 eth-trunk 1
quit
interface GigabitEthernet 0/0/2
 eth-trunk 1
quit
\end{lstlisting}

Configure the matching bundle on the far switch, then verify:

\begin{lstlisting}[style=dcnterminal]
<SW1> display eth-trunk 1
\end{lstlisting}

\medskip\noindent\textbf{Part D --- Why manual mode is dangerous.}\par
Manual mode forwards on every member port unconditionally. It never checks that the far end agrees.

\begin{keyidea}
Miscable one member --- into the wrong switch, or the wrong port --- and manual mode forwards over it anyway. The result is a forwarding loop that spanning tree cannot see, because the trunk hides its members from STP.

That is a broadcast storm, and it takes the segment down. It is the single best argument for LACP.
\end{keyidea}

\medskip\noindent\textbf{Part E --- Static LACP mode.}\par

\begin{lstlisting}[style=dcnterminal]
interface Eth-Trunk 1
 undo mode manual load-balance
 mode lacp-static
quit
\end{lstlisting}

Now each end sends LACP protocol frames announcing its system ID, port priorities, and key. A member joins the bundle only when both ends agree it should. Miscabling produces a port that stays out of the bundle --- a fault that is visible, contained, and safe.

\medskip\noindent\textbf{Part F --- Priorities decide the active links.}\par
LACP allows more member ports than active links. The rest wait as backups.

\begin{lstlisting}[style=dcnterminal]
lacp priority 100                    # system priority; lower wins
interface Eth-Trunk 1
 max active-linknumber 2
quit

interface GigabitEthernet 0/0/1
 lacp priority 100                   # port priority; lower is preferred
quit
\end{lstlisting}

\begin{keyidea}
Two priorities, two different jobs.

\textbf{System priority} decides which switch is in charge of selecting active links. Lower wins; ties break on system MAC. Only one end gets to choose, otherwise the two could disagree.

\textbf{Port priority} is then used by that switch to rank candidate ports. Lower is preferred.

Setting port priorities on the switch that lost the system-priority election has no effect --- a result students often mistake for a broken configuration.
\end{keyidea}

Verify which end is deciding, and which ports are selected:

\begin{lstlisting}[style=dcnterminal]
<SW1> display eth-trunk 1
<SW1> display lacp statistics eth-trunk 1
\end{lstlisting}

\begin{troubleshoot}
\textbf{Trunk stays down.} Member ports differ in speed or duplex. All members must match.

\textbf{Members configured but not selected.} Mode mismatch --- one end manual, the other LACP. Both ends must run the same mode.

\textbf{One port never becomes active.} Expected if \cmd{max active-linknumber} is lower than the member count. Confirm it is a backup, not a fault.

\textbf{Port refuses to join.} It probably still carries per-interface configuration. Ports must be cleared before being added to a trunk.
\end{troubleshoot}

\begin{tryit}
With the trunk up, shut down one active member and watch \cmd{display eth-trunk 1}. Does a backup take over, and how long does it take?

Then, on the switch that lost the system-priority election, change a port priority and confirm nothing happens. Explain why in one sentence.
\end{tryit}
\end{labproject}

\begin{verifybox}
\begin{itemize}
  \item Both manual and LACP modes configured and captured separately.
  \item Both ends use the same mode.
  \item \cmd{display eth-trunk 1} shows members Selected, with the count matching \cmd{max active-linknumber}.
  \item The deciding switch is identified from system priority.
  \item Failover to a backup member is demonstrated and timed.
\end{itemize}
\end{verifybox}

\begin{labdeliverable}
Submit both configurations, \cmd{display eth-trunk} output for each mode, evidence of member failover with timing, and a short note explaining why LACP is preferred to manual mode --- referring specifically to the miscabling scenario.
\end{labdeliverable}
